\begin{recipe}{Tomatensoep van geroosterde tomaten}
\kicker[Voorgerecht,Hartig,Italiaans]{Voorgerecht · hartig · Italiaans}
\index[register]{Tomatensoep van geroosterde tomaten}
\index[register]{Tomaat}
\index[register]{Kippenbouillon}

\meta{50 minuten \dvd Oven 220\,°C}

\lettrine{G}{eroosterde} tomaten geven deze soep een diepe, geconcentreerde smaak.
De tomaten karamelliseren in de oven, en daarna heb je weinig extra nodig.

\blockrule{Ingrediënten}
\begin{ingredients}
  \ing{1 kg}{pruimtomaten, gewassen en gehalveerd}\index[register]{Tomaat}
  \ing{4 el}{olijfolie}
  \ing{1}{ui, in dunne ringen}
  \ing{2 grote}{knoflookteentjes, gehalveerd}
  \ing{1 liter}{kippenbouillon}
  \ing{285 g}{zongedroogde tomaten op olijfolie, uitgelekt}
  \ingb{kleine handvol tijmtakjes}
  \ingb{kleine handvol basilicumtakjes}
  \ingb{versgemalen peper en zout}
  \ingb{kleine basilicumblaadjes voor garnering}
\end{ingredients}

\blockrule{Bereiding}
\tip{Wil je het kerstdinerwaardig, wrijf de soep door een fijne zeef met de achterkant van een lepel is na pureren.}
\begin{steps}
  \item Verwarm de oven voor op 220\,°C (boven- en onderwarmte).
  \item Schenk de olijfolie in een ovenschaal. Schep de tomaten, uienringen en
        knoflook erin en schep goed om. Strooi de tijmtakjes, peper en zout erover.
        Rooster 20--25 minuten
        onder één of twee keer omscheppen, tot alles mooi gekaramelliseerd is. Voeg de
        basilicumtakjes toe in de laatste paar minuten.
  \item Breng de bouillon in een soeppan aan de kook. Schep de geroosterde groenten erin
        (verwijder de tijmtakjes) en voeg de zongedroogde tomaten toe. Laat 5 minuten
        koken.
  \item Pureer het mengsel in de pan glad met een staafmixer. Breng op smaak met peper en zout.
  \item Warm de soep op en schep in verwarmde kommen. Druppel wat olijfolie en leg er wat basilicumblaadjes over.
\end{steps}
\end{recipe}
